\section{Method details}
\label{app:method}

\subsection{Algorithm}
\label{app:algorithm}

SimADEPT runs in three stages. Stage~1 happens entirely in the source domain and
is paid for once. Stages~2 and 3 are the adaptation proper and are the only
parts that consume the target budget of Def.~\ref{def:budget}.

\paragraph{Stage 1: distil the source into a world model (source domain, no
target data).}
We first train an IBRL agent to convergence in the source simulator, which gives
us an expert policy $\pi_s$ and its critic $Q_s$. We then roll out that expert
together with a set of weaker checkpoints saved along its training curve, and
collect a large replay of source trajectories. Using both the expert and the
weaker policies matters: a dataset collected by a single converged policy covers
a narrow tube of states, and the world model is later asked about states outside
that tube.

We then distil this replay into a latent world model with four heads --- an
encoder, a dynamics head, a reward head and a value head --- plus a policy head
that clones the expert's actions. All of them are trained jointly with the
multi-step loss of Eq.~\ref{eq:simdist-distill}: the model encodes a state,
rolls its dynamics forward $H$ steps in latent space, and at each step is asked
to predict the next latent, the reward, the value and the action. The value
targets are relabelled by the frozen source critic, so the value head learns
$V^{\pi_s}$ rather than a return estimate from sparse rewards.

Two details of this stage are load-bearing. The dynamics head is an
\emph{ensemble} of $K$ heads that share the encoder but are independently
initialised and trained on independent bootstrap subsets of each batch, so they
disagree off the data manifold. And the gradient of the policy-cloning term is
stopped before it reaches the encoder, so the policy cannot reshape the latent
space that the dynamics depends on.

\paragraph{Stage 2: warm up the proposal on target demonstrations, tethered to
the source.}
Stage~2 is offline: it uses the $n \le 20$ target demonstrations and no target
interaction. We fine-tune only the policy head, minimising the cloning loss on
the demonstrations plus the tether penalty of Eq.~\ref{eq:tether}. The
encoder, the dynamics, the reward and the value heads are all frozen.

The tether is what makes this stage safe. The demonstrations are few enough that
an unregularised fit revisits them for hundreds of effective epochs, which pulls
the policy far away from the states where the frozen dynamics and value are
still accurate. The tether pool is a fixed set of source states with
$\pi_s$'s actions precomputed on them, so the penalty costs no target
interaction and adds no wall-clock beyond one extra forward pass per step.

\paragraph{Stage 3: repair the dynamics on the agent's own distribution (target
domain, the budget).}
Stage~3 alternates collection and dynamics fitting in rounds. In each round the
agent acts in the target environment with the MPPI planner for a fixed number of
transitions, appends them to a replay buffer, and then takes a number of
gradient steps on the dynamics heads only --- using the same latent-consistency
loss as stage~1, restricted to the dynamics term. The encoder, the reward head,
the value head and the policy head all stay frozen for the entire stage.

Collection uses an optimistic variant of the planner: the return is augmented
with a bonus proportional to the cross-head disagreement, so the agent steers
toward states where its ensemble is uncertain and therefore where the fit has
the most to gain. Evaluation always uses the plain planner with no such bonus.

Freezing the value here is a decision, not an omission: Sec.~\ref{sec:planner-adapt}
shows that adapting it degrades the system in every update order we tried.
Freezing the policy head is what makes the loop on-policy in the sense that
matters --- the states the dynamics is fitted on are exactly the states the
deployed controller will visit, because the controller that collected them is
the controller that will be evaluated.

\begin{algorithm}[h]
\caption{SimADEPT}
\label{alg:simoaed}
\begin{algorithmic}[1]
\Require source simulator $\mathcal{M}_s$, target simulator $\mathcal{M}_t$,
         target demonstrations $\mathcal{D}_t$, budget $N$
\Statex \textbf{Stage 1 --- source distillation}
\State train an IBRL expert $(\pi_s, Q_s)$ in $\mathcal{M}_s$
\State collect source replay $\mathcal{D}_s$ with $\pi_s$ \emph{and} weaker checkpoints of it
\State distil $\{E_\theta, \hat p_\theta^{1..K}, \hat r_\theta, \hat V_\theta, \pi_\theta\}$
       on $\mathcal{D}_s$ with Eq.~\ref{eq:simdist-distill}; relabel values with $Q_s$
\State build tether pool $\mathcal{B} = \{(s, \pi_s(s))\}$ from $\mathcal{D}_s$, frozen
\Statex \textbf{Stage 2 --- proposal warm-up (offline, no target interaction)}
\State fine-tune $\pi_\theta$ only, on $\mathcal{D}_t$ with the tether on $\mathcal{B}$ (Eq.~\ref{eq:tether})
\Statex \textbf{Stage 3 --- dynamics repair (spends the budget)}
\State $\mathcal{R} \gets \mathcal{D}_t$; \; $t \gets 0$
\While{$t < N$}
  \State collect $m$ transitions in $\mathcal{M}_t$ with MPPI, ensemble-mean rollout,
         optimistic disagreement bonus; append to $\mathcal{R}$; \; $t \gets t + m$
  \State take $u$ gradient steps on $\hat p_\theta^{1..K}$ only, latent-consistency
         loss on $\mathcal{R}$ \Comment{$E_\theta, \hat r_\theta, \hat V_\theta, \pi_\theta$ frozen}
\EndWhile
\State \Return planner using $\{E_\theta, \hat p_\theta^{1..K}, \hat r_\theta, \hat V_\theta, \pi_\theta\}$, no bonus
\end{algorithmic}
\end{algorithm}

\subsection{Model structure}
\label{app:model}

Every component is a small MLP; there are no convolutions and no recurrence,
because all of our tasks are state-based.

The \textbf{encoder} maps the raw state to a $512$-dimensional latent. Its output
passes through SimNorm~\citep{hansen2024tdmpc2}, which splits the latent into
groups of $8$ and applies a softmax within each group. This keeps every latent on
a fixed simplex, which is what prevents the magnitude drift that otherwise
accumulates over a multi-step rollout.

The \textbf{dynamics} is an ensemble of $K$ heads, each taking the latent and the
action and predicting the next latent, again SimNorm'd. The heads are separate
MLPs over a shared encoder. At planning time the rollout advances on their mean;
the spread across heads is used only as the collection bonus of stage~3.

The \textbf{reward} and \textbf{value} heads are MLPs on the latent (and the
action, for reward). Both predict a distribution over $101$ bins spanning
$[-10, 10]$ rather than a scalar, and are trained with a two-hot cross-entropy.
The scalar used by the planner is the expectation under that distribution.

The \textbf{policy head} maps the latent to an action. It shares the encoder with
the dynamics, but the cloning gradient is stopped before the encoder, so it reads
the latent without shaping it.

A small \textbf{decoder} maps a latent back to a state. It is trained on
detached encodings only, so it inverts the encoder without influencing it. The
planner does not need it when the policy shares the encoder, which is our default
configuration.

\subsection{Hyperparameters}
\label{app:hparams}

Table~\ref{tab:app-hparams} lists the settings. They are shared across all four
manipulation tasks except for the planner's search parameters, which we re-tune
per task --- we found them not to transfer, including between a learned and an
oracle dynamics model, and we report the swept values in
Table~\ref{tab:app-planner}.

\begin{table}[h]
\centering
\caption{SimADEPT hyperparameters. Shared across tasks unless noted.}
\label{tab:app-hparams}
\small
\begin{tabular}{llc}
\toprule
Stage & Parameter & Value \\
\midrule
\multirow{12}{*}{1 (distillation)}
 & latent dimension & $512$ \\
 & SimNorm group size & $8$ \\
 & MLP hidden size / layers & $512$ / $2$ \\
 & rollout horizon $H$ & $5$ \\
 & per-step loss discount $\rho$ & $0.5$ \\
 & dynamics ensemble size $K$ & $5$ \\
 & reward/value bins, range & $101$, $[-10, 10]$ \\
 & loss weights (dyn.\ / rew.\ / val.\ / BC / dec.) & $20$ / $0.1$ / $0.1$ / $10$ / $1$ \\
 & learning rate & $3\times10^{-4}$ \\
 & batch size & $256$ \\
 & gradient steps & $800$k \\
 & source replay size & $\sim$$55$k episodes \\
\midrule
\multirow{4}{*}{2 (proposal warm-up)}
 & target demonstrations $n$ & $20$ \\
 & gradient steps & $1000$ \\
 & learning rate, batch size & $3\times10^{-4}$, $256$ \\
 & tether strength $\lambda$ & $0.1$--$1.0$ (task-dependent) \\
 & tether pool & $40$k states from $2000$ source episodes \\
\midrule
\multirow{6}{*}{3 (dynamics repair)}
 & target budget $N$ & $100$k transitions \\
 & transitions per round $m$ & $1000$ \\
 & dynamics steps per round $u$ & $1000$ \\
 & learning rate, batch size & $3\times10^{-4}$, $256$ \\
 & collection disagreement bonus & $0.5$ \\
 & adapted components & dynamics heads only \\
\bottomrule
\end{tabular}
\end{table}

\begin{table}[h]
\centering
\caption{MPPI search parameters. \texttt{init\_std} and the base-policy noise
were re-swept per task; the remaining entries are the same everywhere.}
\label{tab:app-planner}
\small
\begin{tabular}{lc}
\toprule
Parameter & Value \\
\midrule
candidate sequences per iteration & $350$ ($100$ from the proposal, $250$ random) \\
planning horizon & $5$ \\
CEM iterations & $3$ \\
elites & $64$ \\
softmax temperature & $0.4$ \\
discount & $0.99$ \\
initial std / minimum std & $0.15$--$1.0$ (swept) / $0.05$ \\
base-policy noise std & $0$--$0.1$ (swept) \\
\bottomrule
\end{tabular}
\end{table}

% ---------------------------------------------------------------------------
% PROVENANCE: every number above was read out of a cfg.yaml of a run that is
% actually reported, not from a dataclass default:
%   stage 1  exps/mimicgen/threading_simdist_dynens_diversev2_h5_seed1/cfg.yaml
%   stage 2+3 exps/square_rotw1_6arm/sq_dynens_anch_d400000_seed2/cfg.yaml
%
% TODO(check): lambda is 0.1 in the Square run above but ~1.0 in the Coffee and
% Threading recipes. Fill the per-task column before submission -- the table
% currently hides a real per-task difference behind a range.
% TODO(check): init_std is 0.15 in the Square adapt run (planner_set override)
% but 1.0 in the "original MPPI" recipe the adapt script applies first. State
% the FINAL value per task; the override order (original_mppi, then planner_set)
% is a known source of confusion.
% TODO(check): stage-1 gradient steps (800k) and source replay size are from the
% Threading cfg. Verify they are the same for the other three tasks; if not,
% split the row per task rather than quoting one.
% TODO(algo): Alg.~\ref{alg:simoaed} needs \usepackage{algorithm} +
% \usepackage{algpseudocode} in the preamble.
% ---------------------------------------------------------------------------
